\documentclass[11pt]{article}
\usepackage[a4paper,margin=25mm]{geometry}
\usepackage{amsmath,amssymb,amsthm}
\usepackage[colorlinks=true,urlcolor=blue]{hyperref}
\newtheorem{theorem}{Theorem}
\newtheorem{definition}[theorem]{Definition}
\newtheorem{remark}[theorem]{Remark}
\title{Adjoining an Additive Absorber That Is a Multiplicative Identity}
\author{}
\date{19 September 2026}
\begin{document}
\maketitle

\paragraph{The question.}
The programme's originator proposes extending the original method of assigning algebraic roles to an adjoined element: after an additive identity/multiplicative absorber $\mathcal T$ and an absorber for both operations $\Omega$, consider an element $u$ satisfying
\begin{equation}\label{roles}
u+x=u=x+u,\qquad ux=x=xu.
\end{equation}
These rules specify a multiplicative identity that is also an additive absorber. The question here is what they preserve when applied to the unchanged operations of the original arithmetic. We give exact constructions and universal receivers. No change in the interpretation of either operation is concealed.

\paragraph{Algebraic conventions.}
An object $D$ has a commutative additive monoid $(D,+,0_D)$, a commutative multiplicative monoid $(D,\cdot,1_D)$ and distributivity. We explicitly state when $0_Dx=0_D$ is required. An operation-preserving, zero-preserving map need not preserve the multiplicative identity unless said to do so. Fix a nonzero commutative unital ring $R$ and the original construction
\[
S_R=R\sqcup\{\mathcal T\},\qquad e=0_R\ne\mathcal T.
\]
The operations on $R$ are unchanged, $\mathcal T$ is the additive identity of $S_R$, and $\mathcal T x=\mathcal T$ for every $x\in S_R$. These are the precursor's original rules \cite{origin}.

\begin{theorem}[An exact existing realization]
The two-element Boolean algebra with addition OR and multiplication AND satisfies \eqref{roles} with $u=1$. More generally, the three-element chain
\[
C_3=\{t<c<u\},\qquad x+y=\max(x,y),\quad xy=\min(x,y)
\]
is a commutative unital standard semiring satisfying \eqref{roles}. Its zero is $t$ and its unit is $u$.
\end{theorem}
\begin{proof}
Minimum and maximum are associative and commutative on a total order. The least element $t$ is the identity for maximum and an absorber for minimum. The greatest element $u$ is the identity for minimum and an absorber for maximum. For distributivity, assume $y\le z$, which covers every pair after interchange. Then
\[
\min(x,\max(y,z))=\min(x,z),
\]
and $\min(x,y)\le\min(x,z)$, so
\[
\max(\min(x,y),\min(x,z))=\min(x,z).
\]
Commutativity supplies the other distributive law. The same argument on the two-element chain gives Boolean OR and AND. Thus the claimed rules hold with all semiring axioms checked.
\end{proof}

\begin{theorem}[Criterion for a genuinely fresh identity]
Let $S$ be a commutative unital standard semiring. On $S^u=S\sqcup\{u\}$ keep the operations of $S$, make $u$ absorbing for addition, and make $u$ the identity for multiplication. These prescriptions produce a standard semiring if and only if
\begin{equation}\label{criterion}
a+ab=a\qquad(a,b\in S).
\end{equation}
In this construction the old $1_S$ remains the identity on the subalgebra $S$, while $u$ is the identity of $S^u$; the inclusion is zero-preserving and operation-preserving, and does not preserve the global multiplicative identity. Condition \eqref{criterion} is equivalent to $1_S+b=1_S$ for all $b\in S$.
\end{theorem}
\begin{proof}
If distributivity holds on $S^u$, then for $a,b\in S$
\[
a=a u=a(u+b)=a u+ab=a+ab,
\]
which proves necessity. Conversely, commutativity holds by the prescribed symmetric rules. An additive triple containing $u$ evaluates to $u$ under either association; all other additive triples use associativity in $S$. A multiplicative triple containing $u$ reduces, in either association, to the product of the remaining arguments; all other triples use associativity in $S$. The additive identity remains $0_S$, since $0_S+u=u$, and $0_Su=0_S$, so it remains a multiplicative absorber. The multiplicative identity is $u$.

For distributivity, triples in $S$ use its old law. If the multiplier is $u$, the equation $u(b+c)=ub+uc$ is an identity. If the multiplier is $a\in S$ and one of the other two arguments is $u$, the equation is precisely $a=a+ab$, established by \eqref{criterion}; if both are $u$, it is $a=a+a$, which follows by taking $b=1_S$ in \eqref{criterion}. This exhausts the cases.

Taking $a=1_S$ in \eqref{criterion} gives $1_S+b=1_S$. Conversely, multiplying that equation by $a$ and using distributivity gives \eqref{criterion}. The inclusion and its identity convention follow from the construction; its failure to be unital is exactly $1_S\ne u$.
\end{proof}

\begin{theorem}[Universal receiver for the original arithmetic]
Let $D$ be an object under the stated conventions with
\[
1_D+d=1_D\qquad(d\in D),
\]
and let $f:S_R\to D$ preserve both operations and the additive identity, without requiring it to preserve $1_R$. Then $f(r)=f(e)$ for every $r\in R$.

There is a universal such receiver, namely the three-element semiring $C_3$ above, with map
\[
i:S_R\longrightarrow C_3,\qquad
i(\mathcal T)=t,\quad i(r)=c\quad(r\in R).
\]
Explicitly, every pair $(D,f)$ as above admits exactly one unital, zero-preserving morphism $F:C_3\to D$ with $F\circ i=f$:
\[
F(t)=0_D,\qquad F(c)=f(e)=f(1_R),\qquad F(u)=1_D.
\]
Thus the universal extension keeps the added role and the distinction between $\mathcal T$ and the supported ring component, but identifies all original ring amplitudes.
\end{theorem}
\begin{proof}
Put $a=f(1_R)$. It is multiplicatively idempotent. Since $1_D+a=1_D$, distributivity gives
\[
a=a1_D=a(1_D+a)=a+a^2=a+a.
\]
Applying preservation of addition, this says $f(1_R+1_R)=f(1_R)$. Add $f(-1_R)$ to both sides. Preservation and associativity give
\[
f(1_R)=f(e).
\]
For every $r\in R$,
\[
f(r)=f(r1_R)=f(r)f(1_R)=f(r)f(e)=f(re)=f(e).
\]
This proves the forced identification without assuming injectivity or any target cancellation.

In fact the target role equations already force additive idempotence and the standard zero law globally: for any $d\in D$, $d=d(1_D+1_D)=d+d$, and $0_D=0_D(1_D+d)=0_D+0_Dd=0_Dd$. Neither law had to be assumed separately.

The map $i$ preserves operations: two ring arguments stay in $R$ under either operation and map to $c$, which satisfies $c+c=c^2=c$; pairs involving $\mathcal T$ follow the zero row of the chain table. The distinguished identities of $S_R$ map as specified.

Let $c_D=f(e)$. The preceding proof gives $c_D=f(1_R)$, and hence $c_D+c_D=c_D$ and $c_D^2=c_D$. Also $0_Dc_D=f(\mathcal T e)=f(\mathcal T)=0_D$, $0_D^2=f(\mathcal T^2)=0_D$, and $0_D1_D=0_D$ by the target's multiplicative identity. Addition with $0_D$ is its identity law, and addition with $1_D$ returns $1_D$ by hypothesis. Multiplication with $1_D$ is its identity law. These checks cover the complete addition and multiplication tables of $\{0_D,c_D,1_D\}$, even if some elements coincide. Therefore the displayed $F$ is a morphism. It is unique because $t,c$ are in the image of $i$ and $u$ must map to the target unit. This proves the universal property.
\end{proof}

\begin{theorem}[The exact unital specialization]
If $f$ in the preceding theorem must also preserve the old multiplicative identity, then the universal receiver is the Boolean semiring and the map is
\[
\chi_R:S_R\longrightarrow\mathbb B,\qquad
\chi_R(\mathcal T)=0,\quad \chi_R(r)=1\quad(r\in R).
\]
It is obtained from $C_3$ by the exact quotient $c\sim u$.
\end{theorem}
\begin{proof}
The preceding proof forces $f(r)=f(1_R)=1_D$ for every ring element. The image consists of $0_D$ and $1_D$, with the Boolean operations: $1_D+1_D=1_D$, $1_D^2=1_D$, and $0_D1_D=0_D$. The unique map from $\mathbb B$ selecting those two elements therefore factors $f$. The quotient identifying $c$ and $u$ in the chain has precisely this two-element table. It does not identify $t$ with them, since $\chi_R$ itself realizes the two distinct classes.
\end{proof}

\paragraph{Consequences for the proposed datum.}
The role assignment exists and has an exact algebraic realization. For the ring-based starting object, adjoining that role has a sharply determined effect: the universal three-state receiver collapses the amplitude within $R$, whereas it retains the support distinction and the newly adjoined unit. Preserving the original unit gives exactly the Boolean support quotient. The distinction between these two receivers is the distinction between a new global identity and the already existing identity of the ring component. These results answer what the proposed extension maintains, rather than deciding its usefulness from the familiarity of a symbol.

\begin{thebibliography}{9}
\bibitem{origin} HI-AI, \emph{An Algebraic Structure Incorporating a Z/1Z-Symmetric Element Adjoined to the Integers: Construction, Analysis, and Generalizations}, recorded publication date 12 July 2025, \href{https://zenodo.org/records/17555345}{Zenodo 17555345}, source \texttt{11.tex}, Section 3. The present $u$ role is the programme originator's explicit proposal in the discussion of 19 September 2026; the proof here establishes its exact receiver.
\end{thebibliography}
\end{document}
